\section{Hiện thực hệ thống}

\subsection{Kiến trúc Đa nhiệm FreeRTOS trên Yolo UNO}
Hệ thống được thiết kế chạy trên một bo mạch Yolo UNO duy nhất với chip xử lý kép ESP32-S3. Để phân phối tài nguyên và chạy song song các khối chức năng, toàn bộ mã nguồn được hiện thực dưới dạng các tác vụ (Tasks) của hệ điều hành thời gian thực FreeRTOS. Tại hàm khởi tạo \texttt{setup()}, các tài nguyên đồng bộ bao gồm Mutex và các Semaphore nhị phân được khởi tạo trước khi phân phối vào từng tác vụ:

\begin{lstlisting}[language=C++]
// Initialize SystemData resources
systemData.dht20 = &dht20_inst;
systemData.lcd = &lcd_inst;
systemData.pixel = &pixel_inst;

systemData.xMutex = xSemaphoreCreateMutex();
systemData.xLedSemaphore = xSemaphoreCreateBinary();
systemData.xNeoSemaphore = xSemaphoreCreateBinary();
xSerialMutex = xSemaphoreCreateMutex();

xTaskCreate(led_blinky, "Task LED Blink", 4096, &systemData, 2, NULL);
xTaskCreate(neo_blinky, "Task NEO Blink", 4096, &systemData, 2, NULL);
xTaskCreate(temp_humi_monitor, "Task TEMP HUMI Monitor", 4096, &systemData, 2, NULL);
xTaskCreate(tiny_ml_task, "Tiny ML Task", 8192, &systemData, 2, NULL);
\end{lstlisting}

Để loại bỏ hoàn toàn các biến toàn cục không an toàn, tất cả các tác vụ đều nhận dữ liệu và tài nguyên thông qua con trỏ cấu trúc \texttt{SystemData} được truyền vào tham số \texttt{pvParameters}.

\subsection{Hiện thực Khối cảnh báo và hiển thị ngoại vi}

\subsubsection{Tác vụ điều khiển LED đơn (Task 1)}
Tác vụ quản lý đèn LED đơn chỉ thị cảnh báo nhiệt độ (\texttt{led\_blinky}) được ghim trong một vòng lặp vô hạn. Tác vụ duy trì trạng thái chặn (\textit{Blocked State}) nhằm tiết kiệm tài nguyên CPU bằng cách chờ đợi tín hiệu từ tác vụ đo đạc thông qua \texttt{xLedSemaphore}:
\begin{lstlisting}[language=C++]
if (xSemaphoreTake(data->xLedSemaphore, portMAX_DELAY) == pdTRUE) {
    // Safely access temperature data using Mutex
    if (xSemaphoreTake(data->xMutex, portMAX_DELAY) == pdTRUE) {
        float current_temp = data->temperature;
        xSemaphoreGive(data->xMutex);

        // Define 3 blinking behaviors based on temperature
        if (current_temp < 30.0) {
            blink_delay = 2000; 
        } 
        else if (current_temp >= 30.0 && current_temp < 35.0) {
            blink_delay = 1000;
        } 
        else {
            blink_delay = 200;
        }
    }
}
\end{lstlisting}
Sau khi giải phóng Mutex, tác vụ sẽ nhấp nháy đèn LED vật lý (\texttt{LED\_GPIO}) với chu kỳ trễ tương ứng bằng hàm \texttt{vTaskDelay()}.

\subsubsection{Tác vụ điều khiển NeoPixel (Task 2)}
Tương tự như tác vụ LED đơn, việc đồng bộ hóa dải LED RGB NeoPixel (\texttt{neo\_blinky}) được hiện thực bằng cơ chế Semaphore nhị phân nhằm tối ưu hóa phản hồi chu kỳ:
\begin{lstlisting}[language=C++]
if (xSemaphoreTake(data->xNeoSemaphore, portMAX_DELAY) == pdTRUE) {
    float current_humi = 0;
    // Safely access humidity data using Mutex
    if (xSemaphoreTake(data->xMutex, portMAX_DELAY) == pdTRUE) {
        current_humi = data->humidity;
        xSemaphoreGive(data->xMutex);
    }

    // Define 3 humidity levels/colors
    if (current_humi < 40.0) {
        data->pixel->setPixelColor(0, data->pixel->Color(255, 0, 0)); // Red (Dry)
    } 
    else if (current_humi >= 40.0 && current_humi <= 70.0) {
        data->pixel->setPixelColor(0, data->pixel->Color(0, 255, 0)); // Green (Normal)
    } 
    else {
        data->pixel->setPixelColor(0, data->pixel->Color(0, 0, 255)); // Blue (Humid)
    }
    data->pixel->show();
}
\end{lstlisting}

\subsubsection{Tác vụ đo đạc cảm biến và hiển thị LCD (Task 3)}
Tác vụ giám sát nhiệt độ - độ ẩm và màn hình hiển thị (\texttt{temp\_humi\_monitor}) thực hiện đọc dữ liệu từ cảm biến DHT20 thông qua bus I2C (cấu hình chân SDA = 11, SCL = 12).
Để tránh tranh chấp tài nguyên bus I2C hoặc tranh chấp ghi màn hình LCD 1602 với các tiến trình khác, tác vụ sử dụng \texttt{xMutex} để độc chiếm quyền truy cập an toàn. Sau khi thu thập dữ liệu cảm biến thành công, tác vụ sẽ:
\begin{enumerate}
    \item Cập nhật dữ liệu vào cấu trúc dùng chung \texttt{SystemData} dưới sự bảo vệ của \texttt{xMutex}.
    \item Phân tích ngưỡng nhiệt độ và cập nhật trạng thái hiển thị LCD theo 3 mức cảnh báo rõ rệt (\texttt{State: Normal}, \texttt{State: Warning}, hoặc \texttt{State: Critical}).
    \item Giải phóng \texttt{xLedSemaphore} và \texttt{xNeoSemaphore} để đánh thức Task 1 và Task 2 thực thi phản hồi chỉ thị.
    \item Sử dụng \texttt{xSerialMutex} để in log kết quả ra cổng Serial, đồng thời đóng gói dữ liệu cảm biến dạng JSON gửi xuống các client WebSocket và đẩy lên đám mây thông qua giao thức MQTT.
\end{enumerate}

\subsection{Hiện thực Khối Webserver và giao thức WebSocket}
Để quản trị và điều khiển thiết bị tại chỗ qua mạng Wi-Fi nội bộ, vi điều khiển ESP32-S3 khởi chạy một máy chủ web không đồng bộ (\texttt{ESPAsyncWebServer}) lắng nghe cổng 80:
\begin{lstlisting}[language=C++]
server.on("/", HTTP_GET, [](AsyncWebServerRequest *request)
          { request->send(LittleFS, "/index.html", "text/html"); });
server.on("/script.js", HTTP_GET, [](AsyncWebServerRequest *request)
          { request->send(LittleFS, "/script.js", "application/javascript"); });
server.on("/styles.css", HTTP_GET, [](AsyncWebServerRequest *request)
          { request->send(LittleFS, "/styles.css", "text/css"); });
\end{lstlisting}

Giao diện WebUI được lưu trữ trực tiếp trong bộ nhớ Flash của bo mạch nhờ hệ thống tệp tin LittleFS. Thay vì cơ chế HTTP Polling liên tục gây tải cao cho chip, hệ thống sử dụng giao thức song công WebSocket để truyền tải thông tin thời gian thực:
\begin{lstlisting}[language=C++]
AsyncWebSocket ws("/ws");
ws.onEvent(onEvent);
server.addHandler(&ws);
\end{lstlisting}

Khi nhận được gói tin điều khiển từ Client, hàm callback \texttt{onEvent} gọi tiến trình \texttt{handleWebSocketMessage()} để giải mã gói tin JSON:
\begin{itemize}
    \item \textbf{Trang thiết bị (device):} Giải mã cổng GPIO cần điều khiển và trạng thái bật/tắt (ON/OFF) từ Client gửi về để tác động trực tiếp lên chân GPIO phần cứng.
    \item \textbf{Trang cấu hình (setting/wifi\_change):} Tiếp nhận các chuỗi cấu hình bao gồm SSID/Password WiFi và Token/Server/Port CoreIOT để ghi nhận và chuyển qua module lưu trữ tệp tin.
    \item \textbf{Khôi phục cấu hình (reset):} Xóa tệp tin lưu cấu hình và thực hiện khởi động lại hệ thống (\texttt{ESP.restart()}).
\end{itemize}

\subsection{Hiện thực Khối kết nối đám mây CoreIOT và mDNS}
Bên cạnh kết nối cục bộ, tác vụ nền định kỳ thực hiện kết nối với nền tảng đám mây CoreIOT thông qua giao thức MQTT bảo mật và thư viện ThingsBoard:
\begin{lstlisting}[language=C++]
void CORE_IOT_reconnect() {
    if (!tb.connected()) {
        if (!tb.connect(g_CORE_IOT_SERVER.c_str(), g_CORE_IOT_TOKEN.c_str(), g_CORE_IOT_PORT.toInt())) {
            return;
        }
        // Publish hardware MAC address as attribute for device identification
        tb.sendAttributeData("macAddress", WiFi.macAddress().c_str());
        tb.sendAttributeData("localIp", WiFi.localIP().toString().c_str());
        tb.RPC_Subscribe(callbacks.cbegin(), callbacks.cend());
        tb.Shared_Attributes_Subscribe(attributes_callback);
    } else {
        tb.loop();
    }
}
\end{lstlisting}

Khi kết nối được thiết lập, thiết bị lập tức gửi lên máy chủ địa chỉ MAC vật lý mặc định (truy xuất trực tiếp từ eFuse) và địa chỉ IP nội bộ làm các thuộc tính định danh duy nhất (\textit{Attributes}). Dữ liệu cảm biến nhiệt độ, độ ẩm được xuất bản định kỳ thành công dưới dạng dữ liệu viễn trắc (\textit{Telemetry}).
Đồng thời, hệ thống triển khai mDNS cho phép người dùng truy cập giao diện Web Server thông qua tên miền cục bộ thân thiện \url{http://yolound.local} mà không cần ghi nhớ địa chỉ IP động của thiết bị.

\subsection{Hiện thực Khối TinyML và Tự kiểm thử (Startup Test)}
Khối học máy cận biên (\texttt{tinyml.cpp}) chạy mô hình phân loại bất thường môi trường trên Tensor Arena tĩnh nhằm tránh phân mảnh bộ nhớ động:

\begin{lstlisting}[language=C++]
constexpr int kTensorArenaSize = 8 * 1024; // 8KB
alignas(16) static uint8_t tensor_arena[kTensorArenaSize];
\end{lstlisting}

\subsubsection{Kịch bản tự kiểm thử khi khởi động (Startup Test)}
Để đảm bảo tính chính xác của mô hình trước khi đi vào vận hành thực tế, hàm \texttt{evaluate\_accuracy()} được gọi tự động khi khởi động. Hàm này nạp lần lượt 50 mẫu dữ liệu kiểm thử (định nghĩa tại \texttt{test\_dataset.h}) vào con trỏ tensor đầu vào của mô hình, chạy suy luận và tính toán ma trận nhầm lẫn (\textit{Confusion Matrix}) bao gồm các chỉ số Accuracy, Precision và Recall:
\begin{lstlisting}[language=C++]
float accuracy = ((float)correct_predictions / (float)TEST_SAMPLES_COUNT) * 100.0;
float precision = 0.0;
if (true_positives + false_positives > 0) {
    precision = ((float)true_positives / (true_positives + false_positives)) * 100.0;
}
float recall = 0.0;
if (true_positives + false_negatives > 0) {
    recall = ((float)true_positives / (true_positives + false_negatives)) * 100.0;
}
\end{lstlisting}

\subsubsection{Tác vụ suy luận thời gian thực TinyML (Task 5)}
Sau khi tự kiểm thử, tác vụ đi vào vòng lặp \texttt{tiny\_ml\_task} chính thức. Tác vụ định kỳ (5 giây một lần) truy xuất an toàn dữ liệu nhiệt độ và độ ẩm mới nhất từ cảm biến DHT20, chuyển đổi sang kiểu dữ liệu số thực \texttt{float32} và chạy suy luận trực tiếp để xác định Điểm bất thường (\textit{Anomaly Score}):
\begin{lstlisting}[language=C++]
input->data.f[0] = temp;
input->data.f[1] = humi;
TfLiteStatus invoke_status = interpreter->Invoke();
float result = output->data.f[0];
\end{lstlisting}

\subsection{Hiện thực Khối quản lý cấu hình động bằng LittleFS}
Hệ thống sử dụng phân vùng bộ nhớ Flash định dạng hệ thống tệp tin LittleFS để lưu trữ các thông số kết nối cục bộ và đám mây tại đường dẫn \texttt{/info.dat} dưới định dạng JSON:
\begin{lstlisting}[language=C++]
void Load_info_File() {
  File file = LittleFS.open("/info.dat", "r");
  if (!file) return;
  DynamicJsonDocument doc(4096);
  DeserializationError error = deserializeJson(doc, file);
  if (!error) {
    g_WIFI_SSID = doc["WIFI_SSID"].as<String>();
    g_WIFI_PASS = doc["WIFI_PASS"].as<String>();
    g_CORE_IOT_TOKEN = doc["CORE_IOT_TOKEN"].as<String>();
    g_CORE_IOT_SERVER = doc["CORE_IOT_SERVER"].as<String>();
    g_CORE_IOT_PORT = doc["CORE_IOT_PORT"].as<String>();
  }
  file.close();
}
\end{lstlisting}

Khi khởi động:
\begin{itemize}
    \item Nếu file cấu hình không tồn tại hoặc rỗng, hệ thống sẽ tự động kích hoạt chế độ Điểm truy cập (\texttt{AP Mode}) để phát Wi-Fi nội bộ. Người dùng truy cập trang cấu hình trên trình duyệt để nhập SSID/Password WiFi mới và cấu hình token cho đám mây CoreIOT.
    \item Khi nhận được thông tin lưu cấu hình qua kênh WebSocket, hàm \texttt{Save\_info\_File()} sẽ serialize tài liệu JSON và ghi đè vào \texttt{/info.dat}, sau đó gọi lệnh \texttt{ESP.restart()} để chuyển đổi sang chế độ Trạm (\texttt{STA Mode}), tự động liên kết Internet và kết nối đám mây CoreIOT.
\end{itemize}